\textbf{COLUNI 2024 -- Questão 37}\par\medskip Observe a imagem abaixo e leia o texto a seguir:

\noindent
\begin{minipage}[t]{0.24\linewidth}
\begin{center}
\includegraphics[width=0.90\linewidth,trim=0 12pt 0 0,clip]{static/simulados/coluni/figuras/2024/questao-37/figura-1.png}\\
{\small Ailton Krenak durante a Assembleia Nacional Constituinte em 1987}
\end{center}
\end{minipage}\hfill
\begin{minipage}[t]{0.73\linewidth}
A imagem se refere ao protesto do líder indígena Ailton Krenak, durante a Assembleia Nacional Constituinte, em 1987. Ocupando a tribuna do Congresso Nacional, Ailton Krenak, proferiu um intenso discurso. Vestindo um terno branco, ele pintou o rosto com pasta de jenipapo, em sinal de luto contra as tentativas de se impedir a aprovação de propostas em defesa do direito à organização social e à ocupação originária e tradicional de terras pelas populações indígenas. Apesar de não ter absorvido todas as propostas apresentadas, tais como o reconhecimento do caráter pluriétnico do Brasil, a Constituição de 1988, acabou por reconhecer, a partir do Artigo 231, os direitos indígenas, incluindo o direito originário a seus territórios tradicionalmente ocupados, que devem ser demarcados pelo Estado brasileiro. Atualmente, há intenso questionamento de grupos econômicos e políticos em relação ao conteúdo do Artigo 231, especialmente em relação à demarcação dos territórios indígenas, representando risco de retrocesso em relação às garantias constitucionais vigentes.
\end{minipage}

Em relação aos direitos indígenas no Brasil contemporâneo, assinale a afirmativa INCORRETA:
\begin{enumerate}
\item A primeira Constituição Brasileira a reconhecer os indígenas como parte da nação brasileira foi a de 1934,
        determinando que fossem respeitadas as terras “que nelas se achem permanentemente localizados”, não
        podendo ser vendidas.
\item O Serviço de Proteção aos Índios (SPI), criado em 1910, foi a primeira instituição do gênero, tendo como
        objetivo a integração total dos indígenas à nação brasileira, a partir de ações assistencialistas, sem que se
        atentasse para o respeito às demandas específicas dos próprios povos.
\item Durante a ditadura civil-militar (1964 - 1985), os povos indígenas foram violentamente atacados, a partir
        do projeto desenvolvimentista defendido pelo governo de exceção, bem como o exercício de controle
        social a partir da chamada Guarda Rural Indígena (GRIN).
\item O direito indígena à terra, é garantido pela Constituição de 1988, entretanto, não se refere a um direito
        originário, ou seja, anterior à criação do próprio Estado, mesmo que se reconheça que os índios foram os
        primeiros ocupantes do Brasil.
\end{enumerate}
